%=============================================================================
% FINDINGS COMPENDIUM: CLASSICAL ITERATION CYCLE (i13--i31)
%
% Compiled: 23 March 2026
% Source: ITERATION_TRACKER.md + all W4i13--W4i31 files
%
% Scope: Every major finding from the ``classical'' cycle, which runs
% from Wave 4 iteration 13 (the first full 17-agent Opus 4.6 formation
% after the Wave 3 convergence and Wave 4 warm-up i10--i12) through
% iteration 31 (the penultimate iteration of the 20-iteration classical
% programme, before the Quantum Gravity cycle i40+).
%
% For each finding: what was found, which iteration, impact (1--4),
% and whether it corrected a prior result.
%=============================================================================

\documentclass[11pt,a4paper]{article}
\usepackage[utf8]{inputenc}
\usepackage[margin=2.5cm]{geometry}
\usepackage{amsmath,amssymb}
\usepackage{booktabs}
\usepackage{longtable}
\usepackage{enumitem}
\usepackage{xcolor}
\usepackage{tcolorbox}
\usepackage{hyperref}
\usepackage{array}

\definecolor{impact4}{RGB}{200,0,0}
\definecolor{impact3}{RGB}{200,120,0}
\definecolor{impact2}{RGB}{0,128,0}
\definecolor{impact1}{RGB}{80,80,80}

\newcommand{\impfour}{\textcolor{impact4}{\textbf{Impact 4 (Paradigm)}}}
\newcommand{\impthree}{\textcolor{impact3}{\textbf{Impact 3 (Major)}}}
\newcommand{\imptwo}{\textcolor{impact2}{\textbf{Impact 2 (Significant)}}}
\newcommand{\impone}{\textcolor{impact1}{\textbf{Impact 1 (Incremental)}}}
\newcommand{\ee}[1]{\times 10^{#1}}
\newcommand{\DFD}{\textsc{dfd}}
\newcommand{\LCDM}{$\Lambda$CDM}

\title{\textbf{Findings Compendium: Classical Iteration Cycle}\\[6pt]
{\large Iterations 13--31 (Wave 4, Opus 4.6)}\\[4pt]
{\normalsize 20--21 March 2026}}
\author{DFD 17-Agent Research Programme}
\date{Compiled 23 March 2026}

\begin{document}
\maketitle
\tableofcontents
\newpage


%=============================================================================
\section{Cycle Overview}
%=============================================================================

The classical iteration cycle spans i13 (20 March 2026) through i31
(21 March 2026). It begins with the first full 17-agent Opus~4.6
formation after Wave~3 convergence (i1--i9) and the Wave~4 warm-up
(i10--i12). The cycle ends at i31, after which the programme transitions
to the Quantum Gravity cycle (i40+) following the alpha formula
breakthrough.

\begin{itemize}
  \item \textbf{Iterations}: 19 (i13--i31)
  \item \textbf{Agents per iteration}: 17 (i13--i23), reduced to 4--6
    focused agents (i24--i31 for some iterations)
  \item \textbf{Total corrections}: $\sim$120+
  \item \textbf{Paradigm corrections}: 3 (distance-bias i14, deep-MOND
    $G_{\rm eff}$ i20, $W'(0)=0$ i20)
  \item \textbf{Prediction count}: 19 (i13) $\to$ 34 (i31)
  \item \textbf{Patent portfolio}: 4 ALIVE, 5 DEAD, 6 NICHE (stable
    from i19 onward)
\end{itemize}


%=============================================================================
\section{Paradigm-Level Findings (Impact 4)}
\label{sec:impact4}
%=============================================================================

%-----------------------------------------------------------------------------
\subsection{$W'(0) = 0$ Correction}
\label{sec:wprime0}
%-----------------------------------------------------------------------------

\begin{tcolorbox}[colback=red!5!white, colframe=red!70!black,
  title=\textbf{$W'(0) = 0$ --- Iteration 20}]

\textbf{Finding:} The DFD formalism file (\texttt{section02\_formalism.tex})
states $W'(0) = 1$, but the simple interpolating function
$\mu(x) = x/(1+x)$ requires $W'(0) = 0$. These are mathematically
incompatible. The $c_s = c$ ``crisis'' from i18 may have been an artefact
of wrong normalisation. Author confirmed $W'(0) = 0$ as the intended
normalisation.

\textbf{Iteration:} i20 (flagged), confirmed at i20 by author.

\textbf{Impact:} \impfour. Resolves the entire $c_s$ crisis thread
(i17--i19). Changes the evaluation point $K''(\bar\Delta)$ from
$\Delta_{\rm bg} \gg 1$ to $\bar\Delta = 0$ (Appendix Q subtraction).
With $W'(0) = 0$, the linearised perturbation theory is degenerate
(Agent~4's approach fails), forcing the deep-MOND quasi-static treatment
(Agent~12) as the correct physical regime.

\textbf{Corrected:} i17 Agent~12's use of $\Delta_{\rm bg} \gg 1$;
i18 Agent~4's $c_s = c$ derivation (which used $K''(0) = 1$ under
incorrect normalisation); the entire ``dust-branch crisis'' narrative
of i17--i19.

\end{tcolorbox}

%-----------------------------------------------------------------------------
\subsection{Deep-MOND $G_{\rm eff}$ Rescue}
\label{sec:deep-mond}
%-----------------------------------------------------------------------------

\begin{tcolorbox}[colback=red!5!white, colframe=red!70!black,
  title=\textbf{Deep-MOND $G_{\rm eff}$ Rescue --- Iterations 19--21}]

\textbf{Finding:} Cosmological perturbations are in the deep-MOND regime
($g_N \ll a_0$). The effective gravitational coupling is
$G_{\rm eff} = G_N \times \sqrt{a_0 / g_N} \sim 44\times$ at
$k = 0.1\;\mathrm{Mpc}^{-1}$. Growth follows $\delta \sim a^2$ (twice
the CDM rate). This overturns Agent~4's linearised result
($G_{\rm eff} = G_N$ flat), which fails because $\mu(0) = 0$ makes
the linearised theory degenerate.

\textbf{Iteration:} i19 (Agent~12 proposes, Agent~4 disagrees); i20
(resolved: Agent~12 is correct; confirmed by Hwang \& Noh 2025,
arXiv:2410.10205).

\textbf{Impact:} \impfour. Rescues DFD cosmology from the ``$\psi$
doesn't cluster'' crisis of i17--i18. The $\psi$-field is stiff
($c_s \sim c$) and does not cluster, but baryons grow under the
enormously enhanced $G_{\rm eff}$, compensating for the absence of CDM
gravitational wells.

\textbf{Corrected:} Agent~4's linearised $G_{\rm eff} = G_N$ (i19);
the $P(k)$ deficit of $10^{-5}$ (i18 Agent~4); the ``existentially
threatened'' cosmological sector status of i17--i18.

\end{tcolorbox}

%-----------------------------------------------------------------------------
\subsection{DFD is a Distance-Bias Theory (Not Modified Expansion)}
\label{sec:distance-bias}
%-----------------------------------------------------------------------------

\begin{tcolorbox}[colback=red!5!white, colframe=red!70!black,
  title=\textbf{Distance-Bias Paradigm Correction --- Iteration 14}]

\textbf{Finding:} DFD does NOT modify the Friedmann equation. The
$\psi$-screen biases flux-based distances ($D_L$) but BAO measures
geometric distances that are unbiased at leading order. The physical
metric is $g_{00} = -c^2 e^{-\psi}$ (single power, NOT $e^{-2\psi}$).

\textbf{Consequences:}
\begin{itemize}
  \item $D_H/r_d$ sign-change at $z = 1.78$: \textbf{RETRACTED} (artefact
    of incorrect modified-Friedmann framework).
  \item $D_H/r_d$ is monotonically suppressed by 0.3--1.5\% relative
    to \LCDM.
  \item DFD survives DESI. CPL tension reduced from 3.5--4.2$\sigma$ to
    $\sim$2.5$\sigma$.
  \item The $\psi$-screen IS what observers interpret as dark energy.
  \item Decisive test pivots from BAO to SNe~Ia $\psi$-screen anisotropy
    + $G_{\rm eff}$-based growth predictions.
\end{itemize}

\textbf{Iteration:} i14.

\textbf{Impact:} \impfour. Invalidates the i12 $D_H/r_d$ sign-change
prediction and multiple derivative predictions. Reframes the entire
observational strategy.

\textbf{Corrected:} i12 $D_H/r_d$ sign-change; i12 $D_H/r_d$ amplitude
table (50$\times$ error); i13 time-delay predictions (3 retracted);
DESI tension estimate.

\end{tcolorbox}


%=============================================================================
\section{Major Findings (Impact 3)}
\label{sec:impact3}
%=============================================================================

%-----------------------------------------------------------------------------
\subsection{$\sigma_8 = 0.83$ Result}
\label{sec:sigma8}
%-----------------------------------------------------------------------------

\begin{tcolorbox}[colback=orange!5!white, colframe=orange!70!black,
  title=\textbf{$\sigma_8 = 0.83$ --- Iteration 21}]

\textbf{Finding:} Deep-MOND $\delta \sim a^2$ growth from
$\delta_b(z_{\rm dec}) \sim 3\ee{-5}$, with Silk damping providing a
critical turnover at $k_{\rm Silk} \approx 0.13\;\mathrm{Mpc}^{-1}$
and peak $P(k)$ at $k_{\rm turn} \approx 0.075\;\mathrm{Mpc}^{-1}$
($R \sim 80\;h^{-1}\;\mathrm{Mpc}$), gives
$\sigma_8 = 0.83\;(+0.22/-0.08)$. Remarkably close to the observed
$0.811 \pm 0.006$.

\textbf{Iteration:} i21 (quantified from i20's range of 0.5--2).

\textbf{Impact:} \impthree. Demonstrates that DFD can reproduce the
amplitude of cosmic structure without CDM. Consistent with the $S_8$
tension (lensing surveys find $\sigma_8 \sim 0.76$--$0.79$). Full range
0.75--1.05 depending on precise Silk scale.

\textbf{Corrected:} i20 range of 0.5--2 sharpened to $0.83\;(+0.22/-0.08)$.

\end{tcolorbox}

%-----------------------------------------------------------------------------
\subsection{CMB Third Peak Deficit Discovery}
\label{sec:cmb-third-peak}
%-----------------------------------------------------------------------------

\begin{tcolorbox}[colback=orange!5!white, colframe=orange!70!black,
  title=\textbf{CMB Third Peak 10--61\% Deficit --- Iterations 21--23}]

\textbf{Finding:} Without CDM infall, baryons are radiation-pressure
supported at $z \sim 1100$. The third acoustic peak is suppressed:
$R_3 \equiv C_\ell^{(3)} / C_\ell^{(1)} \sim 0.55$--$0.70$ vs
observed $\sim 1.0$. This represents a 30--45\% deficit (i21),
narrowed to 10--61\% at i23 after recognising that the third-peak
scale is in the \emph{transition} regime ($g_N / a_0 \sim 1.6$, not
deep-MOND), giving $G_{\rm eff}$ enhancement of only $1.6\times$
(not $7\times$).

\textbf{Iteration:} i21 (quantified); i22 (reduced to 15--30\% with
partial $\psi$-baryon loading); i23 (widened to 10--61\% with transition
regime correction).

\textbf{Impact:} \impthree. The highest-priority unresolved tension of
the entire programme. Only the SymBoltz.jl Boltzmann code can resolve
this numerically. Analogous to the challenge faced by Skordis \& Z\l o\'snik
in AeST (solved with a vector field); DFD has no such mechanism.

\textbf{Corrected:} i21 deep-MOND assumption for third-peak scale
$\to$ i23 transition regime ($g_N/a_0 \sim 1.6$); $G_{\rm eff}$
at third peak: $7\times \to 1.6\times$; $f_\psi$ baryon loading:
$0.15$--$0.30 \to \sim 0.08$.

\end{tcolorbox}

%-----------------------------------------------------------------------------
\subsection{$\beta = 0$ Birefringence Theorem}
\label{sec:beta0}
%-----------------------------------------------------------------------------

\begin{tcolorbox}[colback=orange!5!white, colframe=orange!70!black,
  title=\textbf{$\beta = 0$ Cosmic Birefringence Theorem --- Iterations 19--27}]

\textbf{Finding:} The $S^3$ Chern--Simons term does NOT generate a 4D
$\psi F\tilde{F}$ coupling. Three independent no-go arguments establish
$\beta = 0$ as a structural theorem (not merely leading-order).

\textbf{Evolution:}
\begin{itemize}
  \item i16: First stated as prediction P24.
  \item i19: Upgraded from conjecture to theorem-grade (S$^3$ CS no-go proved).
  \item i20: Tension downgraded from 7$\sigma$ to $\sim$2--3$\sigma$
    (Planck systematic uncertainty $\pm 0.28^\circ$ can absorb the signal;
    Planck and ACT disagree by $>3\sigma$).
  \item i25--i27: Further downgraded as SPIDER data are incompatible with
    Planck/ACT signal. Three-experiment disagreement makes instrumental
    systematics the most likely explanation.
  \item i31: Risk elevated again due to PRL 2026 $n\pi$ phase ambiguity
    resolution. True rotation angle may be larger.
\end{itemize}

\textbf{Iteration:} i19 (theorem grade); i20 (tension quantified).

\textbf{Impact:} \impthree. DFD occupies a unique phenomenological corner:
EM--GW birefringence YES, inter-polarisation NO. If confirmed nonzero by
SO (2027--2028), the microsector dies but core DFD phenomenology survives.

\textbf{Corrected:} i19 elevated from conjecture; i20 $\beta$ tension
from $7\sigma \to 2$--$3\sigma$.

\end{tcolorbox}

%-----------------------------------------------------------------------------
\subsection{SQMS Bound Tightening}
\label{sec:sqms}
%-----------------------------------------------------------------------------

\begin{tcolorbox}[colback=orange!5!white, colframe=orange!70!black,
  title=\textbf{SQMS Bound Tightening --- Iterations 18--19}]

\textbf{Finding:} SRF cavity walls are TRANSPARENT to $\psi$ (couples
via gravitational interaction $-(c^2/2)\psi\rho$, not to $u_{\rm EM}$
directly). This means $Q_{\psi,\rm local} = 20$--$150$ (confirmed i18),
and the SQMS bound tightens from $|\lambda - 1| < 1.56\ee{-9}$ to
$|\lambda - 1| < 2.7\ee{-13}$ (5800$\times$ improvement).

\textbf{Subsequent evolution:}
\begin{itemize}
  \item i16: SQMS absolute bound (1.56$\ee{-9}$) declared UNRELIABLE
    (carries unconstrained $\sqrt{Q_\psi}$ factor). Q-ratio \emph{shape}
    measurement remains ROBUST ($Q_\psi$ cancels).
  \item i18: Wall transparency $\to$ $Q_{\psi,\rm local} = 20$--$150$
    $\to$ bound tightens to $2.7\ee{-13}$.
  \item i19: At this tightened bound, P12 cooperativity drops to
    $1.4\ee{-8}$, re-killing P12.
\end{itemize}

\textbf{Iteration:} i18 (wall transparency + tightening); i16 (absolute
bound unreliable).

\textbf{Impact:} \impthree. Reshapes the entire detection programme.
The Q-ratio shape measurement (psi: 0.52 vs BCS: 3.60, $>30\sigma$
separation) becomes the primary experimental observable. Absolute
sensitivity estimates become secondary.

\textbf{Corrected:} i3 SQMS bound $1.56\ee{-9}$; i12 Dark SRF
reinterpretation (already invalidated at i12); P12 multi-cavity
superradiance (killed again at i19).

\end{tcolorbox}

%-----------------------------------------------------------------------------
\subsection{Patent Portfolio Triage}
\label{sec:patent-triage}
%-----------------------------------------------------------------------------

\begin{tcolorbox}[colback=orange!5!white, colframe=orange!70!black,
  title=\textbf{Patent Portfolio Triage --- Iterations 13--19}]

\textbf{Finding:} Progressive triage of the 15-patent portfolio through
repeated audit cycles:

\begin{center}
\begin{tabular}{lcccc}
\toprule
\textbf{Category} & \textbf{i13} & \textbf{i14} & \textbf{i15} & \textbf{i19 (final)} \\
\midrule
ALIVE & 4 & 4 & 4 & 4 \\
NICHE & 8 & 8 & 7 & 6 \\
THEORETICAL & 1 & 1 & 0 & 0 \\
DEAD & 2 & 3 & 4 & 5 \\
\bottomrule
\end{tabular}
\end{center}

Final portfolio (stable from i19):
\begin{itemize}
  \item \textbf{ALIVE (4):} P5 (Cosmology/Timing), P7 (Pulsed Energy
    Storage), P9 (Positioning), P11 (EM Coupling/Detection)
  \item \textbf{DEAD (5):} P1 (Field Drag), P4 (Wireless Power Transfer),
    P8 (Communications), P12 (Propagation Medium), P13 (Timing Networks)
  \item \textbf{NICHE (6):} P2, P3, P6, P10, P14, P15
\end{itemize}

Key deaths:
\begin{itemize}
  \item P4: i14 --- three no-go arguments ($c_\psi \sim 10^{-5}$ m/s,
    no transverse confinement, geometric dilution)
  \item P12: i15 ($c_\psi$ kills coherence length), briefly revived i17
    ($c_\psi = c$ in lab), re-killed i19 (tightened bound)
  \item P13: i17 (CLONETS-DS irrelevant at SQMS bounds; multi-GNSS
    SNR = 0.09)
\end{itemize}

\textbf{Iterations:} i13--i19 (progressive triage).

\textbf{Impact:} \impthree. Focuses resources on the 4 viable patents.

\textbf{Corrected:} Multiple prior patent viability assessments from
Wave~3.

\end{tcolorbox}

%-----------------------------------------------------------------------------
\subsection{Earth MOND Crossover Correction}
\label{sec:earth-crossover}
%-----------------------------------------------------------------------------

\begin{tcolorbox}[colback=orange!5!white, colframe=orange!70!black,
  title=\textbf{Earth MOND Crossover: 13,000 km $\to$ 12.6 AU --- Iteration 13}]

\textbf{Finding:} The altitude at which the gravitational acceleration
from Earth drops below $a_0$ was incorrectly stated as 13,000--22,000 km
(from W3i3 onward). The correct value is 12.6~AU --- a factor of
$\sim 10^5$ error that had propagated through multiple iterations.

\textbf{Iteration:} i13 (corrected).

\textbf{Impact:} \impthree. All Solar System MOND-regime predictions
are affected. Reinforces that Solar System physics is deeply screened
(Solar EFE gives $a_{\rm sun}/a_0 = 5.3\ee{7}$, as noted at i14).

\textbf{Corrected:} W3i3 value of 13,000--22,000 km, which had
propagated uncorrected through i4--i12.

\end{tcolorbox}

%-----------------------------------------------------------------------------
\subsection{$D_L$ Ratio Correction}
\label{sec:dl-ratio}
%-----------------------------------------------------------------------------

\begin{tcolorbox}[colback=orange!5!white, colframe=orange!70!black,
  title=\textbf{$D_L^{\rm EM} / D_L^{\rm GW}$ Ratio: 1.49 $\to$ 1.315 --- Iteration 19}]

\textbf{Finding:} The luminosity distance ratio
$D_L^{\rm EM} / D_L^{\rm GW} = e^{\Delta\psi(z)}$ was quoted as 1.49
at $z = 1$ with no derivation chain. The correct value is $1.315 =
e^{0.274}$. This is a zero-parameter prediction (Prediction~21), the
cleanest DFD test.

\textbf{Iteration:} i19.

\textbf{Impact:} \impthree. Affects the primary GW--EM comparison
prediction. The 1.49 value had no derivation; 1.315 is derived.

\textbf{Corrected:} Unattributed 1.49 value from earlier iterations.

\end{tcolorbox}

%-----------------------------------------------------------------------------
\subsection{MEMS Reach Correction}
\label{sec:mems-reach}
%-----------------------------------------------------------------------------

\begin{tcolorbox}[colback=orange!5!white, colframe=orange!70!black,
  title=\textbf{MEMS Reach: $10^{-17}$--$10^{-19}$ $\to$ $10^{-14}$--$10^{-16}$ --- Iteration 19}]

\textbf{Finding:} DC gradient MEMS sensitivity was quoted as
$|\lambda - 1| \sim 10^{-17}$ to $10^{-19}$. The wrong source--detector
distance was used. Corrected reach: $10^{-14}$ to $10^{-16}$. Still
$1000\times$ beyond the SQMS bound, so the detection programme survives.

\textbf{Iteration:} i19.

\textbf{Impact:} \impthree. Three orders of magnitude correction to a
key detection pathway.

\textbf{Corrected:} i18 MEMS reach estimate.

\end{tcolorbox}


%=============================================================================
\section{Significant Findings (Impact 2)}
\label{sec:impact2}
%=============================================================================

%-----------------------------------------------------------------------------
\subsection{Axiom Reduction: 4 $\to$ 3 Independent + 1 Derived}
\label{sec:axiom-reduction}
%-----------------------------------------------------------------------------

\begin{tcolorbox}[colback=green!5!white, colframe=green!50!black,
  title=\textbf{Axiom 3 Derived from Axioms 1+2 --- Iteration 19}]

\textbf{Finding:} Axiom~3 ($\mu = x/(1+x)$) is derivable from Axioms~1
and~2 plus the $S^3$ composition law. The minimal axiom set is therefore
3 independent axioms + 1 derived theorem.

\textbf{Iteration:} i19 (Agent~11).

\textbf{Impact:} \imptwo. Strengthens the theoretical economy of DFD.
Reduces the axiom count from 4 to 3 independent + 1 derived. At i23,
further reduction to potentially 1 axiom (A2 derivable from A1) was
proposed but not proven.

\textbf{Corrected:} Prior axiom count of 4 independent.

\end{tcolorbox}

%-----------------------------------------------------------------------------
\subsection{$\psi$-Screen Death / Inflation Replacement Death}
\label{sec:psi-screen}
%-----------------------------------------------------------------------------

\begin{tcolorbox}[colback=green!5!white, colframe=green!50!black,
  title=\textbf{$\psi$-Screen Reframing \& Inflation Status --- Iteration 14}]

\textbf{Finding:} The ``$\psi$-screen'' is reframed at i14 as what
observers interpret as dark energy --- it is the distance-bias mechanism,
not a separate ``screen'' entity. The concept of the $\psi$-screen as an
independent cosmological object effectively dies; it merges into the
distance-bias paradigm.

Regarding inflation replacement: DFD does not contain a slow-roll
inflationary mechanism. The Higgs localisation width $\epsilon_H = 0.05$
was briefly confused with the slow-roll parameter (i42 flag, retracted
i43). DFD is silent on inflation; it neither provides nor replaces it.

\textbf{Iteration:} i14 (distance-bias reframing); i42--i43 ($\epsilon_H$
clarification).

\textbf{Impact:} \imptwo. Simplifies the conceptual framework.

\textbf{Corrected:} Prior ``$\psi$-screen'' as separate entity.

\end{tcolorbox}

%-----------------------------------------------------------------------------
\subsection{P12 Re-Killing}
\label{sec:p12}
%-----------------------------------------------------------------------------

\begin{tcolorbox}[colback=green!5!white, colframe=green!50!black,
  title=\textbf{P12 Death $\to$ Revival $\to$ Death --- Iterations 15--19}]

\textbf{Finding:} P12 (Propagation Medium) has a complex history:
\begin{itemize}
  \item i15: DEAD ($c_\psi$ kills coherence length: $\sim 5\ee{-17}$ m
    at 47~GHz).
  \item i17: REVIVED to NICHE ($c_\psi = c$ in lab confirmed by 3
    independent derivations; coherence length restored to $\sim 0.91$ m).
  \item i19: RE-KILLED. At $|\lambda - 1| < 2.7\ee{-13}$ (tightened
    SQMS bound from i18), cooperativity drops to $1.4\ee{-8}$, output
    power $\sim 10^{-38}$ W.
\end{itemize}

\textbf{Iteration:} Final death at i19.

\textbf{Impact:} \imptwo. Removes one patent from the viable portfolio.

\textbf{Corrected:} i17 revival (itself correcting i15 death).

\end{tcolorbox}

%-----------------------------------------------------------------------------
\subsection{mochi\_class Specification and SymBoltz.jl Identification}
\label{sec:boltzmann}
%-----------------------------------------------------------------------------

\begin{tcolorbox}[colback=green!5!white, colframe=green!50!black,
  title=\textbf{Boltzmann Code: mochi\_class $\to$ SymBoltz.jl --- Iterations 14--29}]

\textbf{Finding:} The optimal Boltzmann code strategy evolved
significantly:
\begin{itemize}
  \item i12: hi\_class identified as starting point; 2520 lines, 4.5 months.
  \item i14: mochi\_class supersedes hi\_class; 500--800 lines, 2--3 months,
    \$40--80K.
  \item i15: mochi\_class background module RETRACTED (DFD doesn't modify
    $H(z)$). 4-module design, $\sim$1850 lines.
  \item i17: mochi\_class developer spec COMPLETE (19 functions, 4 modules,
    2 CLASS hooks).
  \item i19: ``Zero new code'' strategy with existing mochi\_class DEAD
    (Horndeski $\alpha_K$ is linearised $\to$ reproduces degenerate
    Agent~4 result). Needs custom nonlinear AQUAL solver: $\sim$250 lines
    C + 80 lines Python. 3--4 developer-weeks, \$15--20K.
  \item i21: mochi\_class developer-ready specification with 180 lines C
    code (exact quadratic $G_{\rm eff}$ solver).
  \item i29: SymBoltz.jl (published A\&A March 2026) identified as
    primary Boltzmann solver. Symbolic equation input, no approximation
    switching, automatic differentiation. 2--3 month DFD implementation.
    Strategic shift: SymBoltz.jl primary, mochi\_class fallback.
\end{itemize}

\textbf{Iteration:} i14 (mochi\_class); i29 (SymBoltz.jl).

\textbf{Impact:} \imptwo. The Boltzmann code is the single rate-limiting
step for the entire programme. All 7 computation-only open questions at
i31 require it.

\textbf{Corrected:} Multiple prior specifications (hi\_class $\to$
mochi\_class $\to$ SymBoltz.jl).

\end{tcolorbox}

%-----------------------------------------------------------------------------
\subsection{Clock Pre-Registration}
\label{sec:clock}
%-----------------------------------------------------------------------------

\begin{tcolorbox}[colback=green!5!white, colframe=green!50!black,
  title=\textbf{Clock Pre-Registration Paper --- Iterations 20--22}]

\textbf{Finding:} Nuclear clock pre-registration paper developed through
multiple iterations:
\begin{itemize}
  \item i20: Honest negative. With screening: amplitudes $10^{-20}$ to
    $10^{-22}$ (undetectable). Without screening: $\sim 10^{-16}$ but
    inconsistent with v3.2. Paper rewritten with both scenarios.
  \item i22: Nuclear clock v3 paper VALIDATED. All amplitude calculations
    confirmed ($10^{-18}$ Earth-surface, $10^{-16}$ solar-orbit,
    $10^{-13}$ unscreened-excluded). 3000:1 ratio derivation confirmed.
  \item QG cycle i40: Th-229 nuclear clock breakthrough (220 Hz
    reproducibility, Nature 2026). $K = 5900 \pm 2300$ confirmed.
\end{itemize}

\textbf{Iteration:} i20 (initial); i22 (validated); QG i40 ($K$ confirmed).

\textbf{Impact:} \imptwo. Establishes the nuclear clock as the
definitive $\dot\alpha/\alpha$ test for DFD. Nearest test: PTB E3/E2.

\textbf{Corrected:} i20 clock amplitudes from $10^{-16}$ to
$10^{-20}$--$10^{-22}$ (with screening).

\end{tcolorbox}

%-----------------------------------------------------------------------------
\subsection{Bullet Cluster JWST Relief}
\label{sec:bullet-cluster}
%-----------------------------------------------------------------------------

\begin{tcolorbox}[colback=green!5!white, colframe=green!50!black,
  title=\textbf{Bullet Cluster: JWST 10:1 Minor Merger --- Iteration 23}]

\textbf{Finding:} JWST Finner et al.\ 2025 establishes the Bullet Cluster
as a 10:1 minor merger (not the previously assumed 3:1 major merger).
The DFD mass deficit drops from $2.3$--$3.5\times$ (i19) to
$1.5$--$2.5\times$ (i23). With $\sim$2~eV neutrinos: $1.2$--$2.0\times$,
within observational systematics. No longer a clear falsification.

\textbf{Iteration:} i23.

\textbf{Impact:} \imptwo. Relieves the most prominent ``honest negative''
in the DFD programme.

\textbf{Corrected:} i19 mass deficit $2.3$--$3.5\times$.

\end{tcolorbox}

%-----------------------------------------------------------------------------
\subsection{Nuclear Clock $K = 5900$}
\label{sec:nuclear-clock}
%-----------------------------------------------------------------------------

\begin{tcolorbox}[colback=green!5!white, colframe=green!50!black,
  title=\textbf{Th-229 Nuclear Clock $K = 5900 \pm 2300$ --- QG Cycle i40}]

\textbf{Finding:} The Th-229 nuclear clock $\alpha$-sensitivity
coefficient is $K = 5900 \pm 2300$ (Banerjee et al.\ 2025, Nature
Communications). Frequency reproducibility of 220~Hz demonstrated
(Nature 2026). DFD predicts $\dot\alpha / \alpha = 0$ (Prediction~P69),
testable at $\sim 10^{-19}$ by 2029--2030 with single-ion clocks.

\textbf{Iteration:} QG i40 (confirmation and integration).

\textbf{Impact:} \imptwo. Provides a concrete experimental test of
the DFD prediction that $\alpha$ does not vary.

\textbf{Corrected:} No prior value corrected; new experimental data
integrated.

\end{tcolorbox}

%-----------------------------------------------------------------------------
\subsection{Wide Binary Resolution}
\label{sec:wide-binaries}
%-----------------------------------------------------------------------------

\begin{tcolorbox}[colback=green!5!white, colframe=green!50!black,
  title=\textbf{Wide Binaries: $+42\% \to +10$--$12\%$ with EFE --- Iteration 11}]

\textbf{Finding:} The DFD wide binary velocity excess was originally
quoted as $+42\%$. With the Milky Way external field effect (EFE), this
reduces to $+10$--$12\%$. Banik et al.\ 19$\sigma$ Newtonian result
does NOT falsify DFD because the EFE suppresses the MOND signal in the
Solar neighbourhood. Cookson et al.\ (2026) definitively support
Newtonian dynamics in wide binaries, which is GREEN for DFD.

\textbf{Iteration:} i11 (EFE correction); i13 (caution flags); QG i40
(Cookson et al.\ GREEN).

\textbf{Impact:} \imptwo. Removes a potential falsification threat.

\textbf{Corrected:} W3i1 wide binary formula ($+42\% \to +10$--$12\%$).

\end{tcolorbox}

%-----------------------------------------------------------------------------
\subsection{DESI DR2 Match}
\label{sec:desi}
%-----------------------------------------------------------------------------

\begin{tcolorbox}[colback=green!5!white, colframe=green!50!black,
  title=\textbf{DESI DR2 Dark Energy Match --- Iterations 14--27}]

\textbf{Finding:} DFD predicts $w_0 = -0.70 \pm 0.12$,
$w_a = -0.85 \pm 0.25$ (Prediction P28). DESI DR2 measures
$w_0 = -0.752 \pm 0.066$, $w_a = -0.86^{+0.24}_{-0.21}$. The offset
is $< 0.5\sigma$. DFD is one of the best-fitting theoretical frameworks
for DESI data. Key discriminator: DFD forbids phantom crossing
($w > -1$ always); DESI best-fit crosses at $z \sim 0.5$.

\textbf{Iteration:} i14 (initial confrontation); i25 (P28 formalised);
i27 ($< 0.5\sigma$ offset confirmed).

\textbf{Impact:} \imptwo. DFD is competitive with generic $w_0 w_a$CDM
and better than $\Lambda$CDM ($3$--$4\sigma$ disfavoured by DESI).

\textbf{Corrected:} i13 DESI tension (3.5--4.2$\sigma \to \sim 2.5\sigma$
at i14; further refined through i27).

\end{tcolorbox}

%-----------------------------------------------------------------------------
\subsection{Euclid Pre-Registration Planning}
\label{sec:euclid}
%-----------------------------------------------------------------------------

\begin{tcolorbox}[colback=green!5!white, colframe=green!50!black,
  title=\textbf{Euclid DR1 Pre-Registration --- Iterations 17--31}]

\textbf{Finding:} Euclid DR1 (21 October 2026) pre-registration paper
progressively developed:
\begin{itemize}
  \item i17: Euclid DR1 preparation initiated. Tomographic $S_8$ gradient
    3--5$\sigma$. Strategic window identified (no team testing Yukawa
    $G_{\rm eff}(k)$).
  \item i23: $f \cdot \sigma_8$ ratio revised down: $1.05$--$1.15$ (not
    $1.1$--$1.8$). $\sim$2--3$\sigma$ across 4 combined bins.
  \item i31: Draft complete. 7 observables ($w_0$, $w_a$, $\mu_0$,
    $\Sigma_0$, $\eta_0$, $E_G$, $\gamma$), 5 falsification criteria.
    Target: arXiv before 24 June 2026.
\end{itemize}

\textbf{Iteration:} i17 (initiated); i31 (draft complete).

\textbf{Impact:} \imptwo. First DFD paper with explicit, pre-registered,
falsifiable predictions. Deadline 77 days from i31.

\textbf{Corrected:} i23 $f \cdot \sigma_8$ ratio ($1.1$--$1.8 \to
1.05$--$1.15$).

\end{tcolorbox}


%=============================================================================
\section{Other Significant Findings (Impact 2)}
\label{sec:other-impact2}
%=============================================================================

%-----------------------------------------------------------------------------
\subsection{$c_\psi$ Resolution}
%-----------------------------------------------------------------------------

\textbf{Iteration:} i15--i17.

\textbf{Finding:} The propagation speed of $\psi$ perturbations caused
a protracted debate:
\begin{itemize}
  \item i15: $c_s^2 \to 0$ theorem kills ALL resonant scalar-wave sensors.
  \item i16: CONTRADICTORY results (Agent~2: $c_\psi \sim 10^{-13}$ m/s;
    Agent~4: $c_\psi = c$ in lab). Contradiction unresolved.
  \item i17: RESOLVED with 3 independent confirmations. $c_\psi = c$ in lab.
    Agent~2's error: $K'' \to 0$ means the correction vanishes, not speed
    suppression. Action normalisation is presentation error ($c^4/4$ factor).
\end{itemize}

\textbf{Impact:} \imptwo. Revives P12 (briefly) and establishes that
laboratory $\psi$ perturbations propagate at $c$.

%-----------------------------------------------------------------------------
\subsection{GW Never Gravitationally Lensed (Structural Theorem)}
%-----------------------------------------------------------------------------

\textbf{Iteration:} i15 (discovery); i16 (elevated to structural theorem).

\textbf{Finding:} GW propagate on flat $\mathbb{R}^3$ with zero
$\psi$-coupling at ALL post-Newtonian orders. A single multiply-lensed
GW event would FALSIFY DFD. EM images displaced 30--120 arcsec from GW
position. This is DFD's Tier~0 falsification test.

\textbf{Impact:} \imptwo. The cleanest structural prediction. Testable
in the 3G detector era ($\sim$2035+).

%-----------------------------------------------------------------------------
\subsection{SNe Ia $\psi$-Screen Anisotropy}
%-----------------------------------------------------------------------------

\textbf{Iteration:} i15.

\textbf{Finding:} $C_\ell \propto \ell^{-2}$, peaks at $\ell \sim 3$--$15$,
RMS 0.8--1.2\% (Prediction~P20). Testable at $\sim 2\sigma$ with
Pantheon+ NOW; $5\sigma$ with Rubin~Y1.

\textbf{Impact:} \imptwo.

%-----------------------------------------------------------------------------
\subsection{Screening Function Gap and Disformal Resolution}
%-----------------------------------------------------------------------------

\textbf{Iteration:} i23 (gap identified); i24 (resolution path); i25
(gap RESOLVED).

\textbf{Finding:} $\Sigma(y) = 1/\sqrt{1+y}$ CANNOT be derived from
$\mu(x) = x/(1+x)$ via k-essence action. Resolution: DFD recast as
disformal scalar-tensor theory (DHOST Class~Ia). Conformal factor sets
$\mu$; disformal factor sets $\Sigma$ independently. Ghost-free,
gradient-stable, subluminal, $c_T = c$.

\textbf{Impact:} \imptwo. Upgrades Axiom~A2 from pure k-essence to
disformal scalar-tensor. Static predictions unchanged; cosmological
predictions enriched.

%-----------------------------------------------------------------------------
\subsection{Dark Energy EoS Prediction (P28)}
%-----------------------------------------------------------------------------

\textbf{Iteration:} i25 (formalised); refined through i27.

\textbf{Finding:} $w_0 = -0.70 \pm 0.12$, $w_a = -0.85 \pm 0.25$,
with $w \geq -1$ always (no phantom crossing). Overlaps DESI DR2 at
$< 0.5\sigma$. Key discriminator: DFD forbids phantom crossing.

\textbf{Impact:} \imptwo.


%=============================================================================
\section{Incremental Findings (Impact 1)}
\label{sec:impact1}
%=============================================================================

\begin{longtable}{p{1cm}p{3cm}p{7cm}p{2cm}}
\toprule
\textbf{Iter.} & \textbf{Finding} & \textbf{Description} & \textbf{Corrected?} \\
\midrule
\endfirsthead

i13 & T4 Cold Spot RESOLVED & $x_{\rm eff} = 0.10 \pm 0.03$; ISW
$= -63\;(+25/-35)\;\mu$K vs observed $-75 \pm 35\;\mu$K; tension
$0.3\sigma$. All five tensions T1--T5 resolved. & Yes (i12 $x_{\rm eff}$) \\

i13 & W4i12 $D_H/r_d$ table 50$\times$ error & Analytic formula gives
20--100\% deviations at $z > 1$; table showed 1--2\%. & Yes \\

i13 & CPL tension 3.5--4.2$\sigma$ & Exceeds DFD's own falsification
threshold under CPL mapping. & Yes (i14 resolves) \\

i13 & Psi-QEC advantage narrowed & $10$--$83\times \to 3$--$83\times$
vs qLDPC (Iceberg Pinnacle). & Yes \\

i13 & Zero scalar GW memory & New Prediction~13 (zero-parameter). PTA
breathing mode kills DFD. & No (new) \\

i13 & Galactic bar pattern speed & New Prediction~16 ($+19\%$, Gaia
DR4). & No (new) \\

i13 & Excess large voids & $+24\%$, testable $4$--$6\sigma$ with DESI
DR2. & No (new) \\

i14 & P4 DEAD & Three no-go arguments. Distance-independence unproven. &
Yes (i13 P4 alive) \\

i14 & Birefringence 39~ns & Full action-principle derivation. Double
pulsar differential Shapiro. & Yes (i13 25~ns) \\

i14 & mochi\_class $\to$ hi\_class & 2520 $\to$ 500--800 lines. &
Yes \\

i14 & Pairwise kSZ velocity & Prediction~17: $+10\%$ at $> 100\;
h^{-1}$ Mpc. & No (new) \\

i15 & 67$^\circ$K blocker RESOLVED & Analytical Lindblad proof. &
Yes \\

i15 & $\omega^2$ loss scaling DERIVED & $R_Q = 0.50$ (consistent with
0.49). & Yes (prior unverified) \\

i15 & $a_0$ is fundamental constant & NOT a free function. Promoting
$a_0(t)$ breaks convexity. & No (clarification) \\

i16 & Q-ratio RESOLVED: 0.49 correct & W4i15 ambiguity was
double-counting. $0.52 \pm 0.08$ (with TESLA geometry). & Yes \\

i16 & P24 birefringence null & $\beta = 0$ exactly. 3 new predictions
(P22 FRB, P23 21cm, P24). & No (new) \\

i16 & Phase~0 pathfinder & \$50K, 2 weeks, single TESLA 9-cell cavity.
Go/no-go for \$3.2M Phase~I. & No (new) \\

i17 & $c_\psi = c$ in lab & 3 independent confirmations. Action Eq.~(6)
normalisation is presentation error. & Yes (i16 contradiction) \\

i17 & P12 revived to NICHE & Coherence length restored to $\sim$0.91~m
at 47~GHz. & Yes (i15 death) \\

i17 & GPS 59ps RETRACTED & Signal $10^6$ below noise. Replaced with LLR
(0.93~ns at APOLLO). & Yes \\

i17 & H$_0$ tension: only 0.4 km/s/Mpc & DFD cannot fully explain
Hubble tension. Downgraded. & Yes (prior overclaim) \\

i17 & Einstein Telescope decisive & $5\sigma$ in 1--3 years of operation
($\sim$2036--2038). & No (new) \\

i18 & $c_s = c$ CONFIRMED & Full derivation from corrected action.
Robust result. & Yes (confirms i17) \\

i18 & SRF walls transparent & $\psi$ couples gravitationally, not to
$u_{\rm EM}$ directly. & Yes (prior unknown) \\

i18 & DC gradient MEMS & $|\lambda - 1| \sim 10^{-17}$ to $10^{-19}$
(later corrected to $10^{-14}$--$10^{-16}$ at i19). & No (new) \\

i19 & $\beta = 0$ theorem-grade & S$^3$ CS no-go proved. 3 independent
arguments. & Yes (i16 conjecture) \\

i19 & Wall transparency hurts P7 & $Q_{\psi,\rm local} = 20$--$150$
means $\psi$-energy radiates in $<1\;\mu$s. & Yes (P7 storage time) \\

i20 & $\beta$ tension $7\sigma \to 2$--$3\sigma$ & Planck systematics
$\pm 0.28^\circ$; Planck/ACT disagree $> 3\sigma$. & Yes \\

i20 & mochi\_class: 250 lines C & ``Zero new code'' DEAD. 3--4
developer-weeks, \$15--20K. & Yes (i19 strategy) \\

i22 & Nuclear clock v3 validated & Amplitudes confirmed ($10^{-18}$,
$10^{-16}$, $10^{-13}$). & No (validation) \\

i23 & Bullet Cluster JWST & 10:1 minor merger. Deficit $2.3$--$3.5\times
\to 1.5$--$2.5\times$. & Yes \\

i23 & Third-peak transition regime & $g_N / a_0 \sim 1.6$ (not
deep-MOND). $G_{\rm eff} = 1.6\times$ (not $7\times$). & Yes (i21) \\

i23 & Euclid $f \cdot \sigma_8$ revised & $1.1$--$1.8 \to 1.05$--$1.15$.
$\sim$2--3$\sigma$. & Yes \\

i23 & New Prediction P27 & BAO amplitude $\propto \sqrt{a_0}$. & No (new) \\

i24 & $\Sigma(y)$ gap identified & k-essence gives $(1+y)^{-1/3}$, not
$1/\sqrt{1+y}$. Disformal resolution path. & Yes (prior assumed) \\

i25 & $\Sigma(y)$ gap RESOLVED & Disformal calculation complete.
Ghost-free, subluminal, Horndeski-class. & Yes (i24 gap) \\

i25 & P28 dark energy EoS & $w_0 = -0.65 \pm 0.15$, $w_a = -0.8 \pm 0.3$.
Formalised. & No (new) \\

i26 & Birefringence threat downgraded & SPIDER incompatible with
Planck/ACT. Threat yellow $\to$ green-yellow. & Yes (i25 tension) \\

i26 & P30, P31 new predictions & Velocity-dependent screening; gravitational
slip $\eta \approx 0.75$. & No (new) \\

i27 & P32--P34 new predictions & Induced GW spectral scaling;
$A_L$--disformal lensing; DM--DE effective phantom crossing. & No (new) \\

i28 & $\sigma_8$ flag & SPT+ACT+Planck combined: $0.8137 \pm 0.0038$.
DFD: $0.83$. Offset $4.3\sigma$ in \LCDM, but GR-assumed. & No (flag) \\

i29 & $\sigma_8$ flag resolved & $<1\sigma$ offset with proper DFD error
bars and modified gravity extraction. & Yes (i28 flag) \\

i29 & SQMS 2.0 funded \$125M & Parasitic search mode path forward. & No (external) \\

i29 & SymBoltz.jl primary & Published A\&A March 2026. 2--3 month DFD
implementation. & Yes (mochi\_class primary) \\

i30 & Euclid $\mu$--$\Sigma$ forecast & DFD predictions fall within
Euclid's projected $3\sigma$ detection region. & No (validation) \\

i31 & DFD resolves $\nu$-mass crisis & Under DFD modified growth,
$\sum m_\nu$ bound relaxes to $0.10$--$0.16$ eV (from $< 0.064$ eV
under GR). & No (new) \\

i31 & Euclid pre-reg draft complete & 7 observables, 5 falsification
criteria. First pre-registered DFD predictions. & No (deliverable) \\

i31 & SQMS white paper final & Parasitic search, \$75--100K. Ready for
SQMS-NOTE formatting. & No (deliverable) \\

i31 & SymBoltz funding proposal & \pounds 65K / 12 months. Self-funded
pilot \pounds 18K / 3 months. & No (deliverable) \\

\bottomrule
\end{longtable}


%=============================================================================
\section{Master Correction Ledger}
\label{sec:corrections}
%=============================================================================

\begin{longtable}{p{1cm}p{5cm}p{5cm}p{2cm}}
\toprule
\textbf{Iter.} & \textbf{What Was Corrected} & \textbf{Old $\to$ New} &
\textbf{Impact} \\
\midrule
\endfirsthead

i13 & Earth MOND crossover & 13,000--22,000 km $\to$ 12.6 AU ($10^5$ error) & 3 \\
i13 & W4i12 $D_H/r_d$ amplitudes & $\sim$50$\times$ too small & 3 \\
i13 & Psi-QEC advantage & $10$--$83\times \to 3$--$83\times$ & 1 \\
i14 & $D_H/r_d$ sign-change & RETRACTED (wrong framework) & 4 \\
i14 & DESI tension & $3.5$--$4.2\sigma \to \sim 2.5\sigma$ & 3 \\
i14 & P4 & ALIVE $\to$ DEAD & 2 \\
i14 & Time-delay predictions & 3 RETRACTED & 2 \\
i14 & Birefringence delay & 25 ns $\to$ 39 ns & 1 \\
i14 & Boltzmann code & hi\_class $\to$ mochi\_class & 2 \\
i15 & P12 & THEORETICAL $\to$ DEAD & 2 \\
i15 & EM--GW separation & $\sim$1 arcsec $\to$ 30--120 arcsec & 2 \\
i15 & Boltzmann $\alpha_M$ mapping & WRONG; correct implementation specified & 2 \\
i16 & Q-ratio ambiguity & $0.49$ vs $0.27$ $\to$ $0.49$ definitive & 2 \\
i16 & SQMS absolute bound & $1.56\ee{-9} \to$ UNRELIABLE & 3 \\
i16 & Dust-branch $w$ & $10^{-37} \to 5\ee{-6}$ & 1 \\
i17 & $c_\psi$ contradiction & Regime-dependent $\to$ $c$ in lab & 3 \\
i17 & Action Eq.~(6) normalisation & Presentation error (RESOLVED) & 2 \\
i17 & Q-ratio & $0.49 \to 0.52 \pm 0.08$ (geometry correction) & 1 \\
i17 & P12 & DEAD $\to$ NICHE (revived) & 2 \\
i17 & H$_0$ tension & $\sim$60--70\% explained $\to$ only $\sim$7\% & 2 \\
i17 & GPS 59ps & RETRACTED & 2 \\
i18 & SRF wall transparency & Unknown $\to$ TRANSPARENT & 3 \\
i18 & $Q_{\psi,\rm local}$ & Uncertain $\to$ 20--150 (confirmed) & 3 \\
i18 & SQMS bound & $1.56\ee{-9} \to 2.7\ee{-13}$ (5800$\times$) & 3 \\
i19 & $D_L$ ratio & $1.49 \to 1.315$ & 3 \\
i19 & MEMS reach & $10^{-17}$--$10^{-19} \to 10^{-14}$--$10^{-16}$ & 3 \\
i19 & P12 & NICHE $\to$ DEAD (re-killed) & 2 \\
i19 & Axiom count & 4 independent $\to$ 3 + 1 derived & 2 \\
i19 & $\beta = 0$ & Conjecture $\to$ theorem-grade & 2 \\
i19 & $\beta$ tension & $4.4\sigma \to 7\sigma$ (updated data) & 1 \\
i20 & $G_{\rm eff}$ cosmological & $G_N$ flat $\to$ deep-MOND $40$--$140\times$ & 4 \\
i20 & $W'(0)$ & $1 \to 0$ (formalism inconsistency resolved) & 4 \\
i20 & $\beta$ tension & $7\sigma \to 2$--$3\sigma$ (systematics) & 2 \\
i20 & $\sigma_8$ & 0.003 (dead) $\to$ 0.5--2 (viable) & 4 \\
i20 & mochi\_class & Parameter file $\to$ 250 lines C & 2 \\
i20 & Clock amplitudes & $10^{-16} \to 10^{-20}$--$10^{-22}$ (screened) & 2 \\
i21 & $\sigma_8$ & $0.5$--$2 \to 0.83\;(+0.22/-0.08)$ & 3 \\
i21 & CMB third peak & ``potential tension'' $\to$ 30--45\% deficit & 3 \\
i22 & CMB third peak & $30$--$45\% \to 15$--$30\%$ deficit & 1 \\
i23 & Third-peak regime & Deep-MOND $\to$ transition ($g_N/a_0 \sim 1.6$) & 2 \\
i23 & $G_{\rm eff}$ at third peak & $7\times \to 1.6\times$ & 2 \\
i23 & Bullet Cluster deficit & $2.3$--$3.5\times \to 1.5$--$2.5\times$ & 2 \\
i23 & Euclid $f \cdot \sigma_8$ ratio & $1.1$--$1.8 \to 1.05$--$1.15$ & 1 \\
i24 & $\Sigma(y)$ derivation gap & k-essence fails; disformal path identified & 2 \\
i25 & $\Sigma(y)$ & GAP RESOLVED (disformal) & 2 \\
i26 & Birefringence threat & Yellow $\to$ green-yellow & 1 \\
i26 & P28 & $w_0 = -0.65 \to -0.70$; $w_a = -0.8 \to -0.85$ & 1 \\
i29 & $\sigma_8$ ``tension'' & $4.3\sigma \to < 1\sigma$ (proper treatment) & 1 \\
i29 & Boltzmann primary & mochi\_class $\to$ SymBoltz.jl & 2 \\

\bottomrule
\end{longtable}


%=============================================================================
\section{Summary Statistics}
\label{sec:statistics}
%=============================================================================

\subsection{Correction Rate by Iteration}

\begin{center}
\begin{tabular}{lccl}
\toprule
\textbf{Iteration} & \textbf{Corrections} & \textbf{Paradigm?} & \textbf{Theme} \\
\midrule
i13 & $\sim$10 & No & Full 17-agent deep dive \\
i14 & $\sim$15 & Yes & Distance-bias paradigm correction \\
i15 & $\sim$8 & No & Post-paradigm stabilisation \\
i16 & $\sim$12 & No & $c_\psi$ contradiction; corpus rewrite \\
i17 & $\sim$10 & No & $c_\psi$ resolved; dust branch flagged \\
i18 & $\sim$6 & No & Dust branch crisis; SQMS tightening \\
i19 & $\sim$12 & No & Deep-MOND rescue attempt; axiom reduction \\
i20 & $\sim$8 & Yes & $W'(0) = 0$; $G_{\rm eff}$ resolution \\
i21 & $\sim$5 & No & $\sigma_8 = 0.83$; third peak quantified \\
i22 & $\sim$3 & No & Clock validated; third peak improved \\
i23 & $\sim$5 & No & Screening gap; Bullet Cluster JWST \\
i24 & $\sim$2 & No & $\Sigma(y)$ resolution path \\
i25 & $\sim$3 & No & $\Sigma(y)$ RESOLVED; P28--P29 \\
i26 & $\sim$2 & No & Birefringence downgraded; P30--P31 \\
i27 & $\sim$2 & No & P32--P34; DESI match refined \\
i28 & $\sim$2 & No & P33 corrected (lensing mechanism) \\
i29 & $\sim$3 & No & $\sigma_8$ flag resolved; SQMS 2.0; SymBoltz \\
i30 & $\sim$1 & No & Synthesis; final consistency matrix \\
i31 & $\sim$1 & No & Penultimate; publication deliverables \\
\midrule
\textbf{Total} & $\sim$\textbf{110} & 3 paradigm & \\
\bottomrule
\end{tabular}
\end{center}

\subsection{Prediction Count Evolution}

\begin{center}
\begin{tabular}{lccc}
\toprule
\textbf{Iteration} & \textbf{Total Predictions} & \textbf{New} & \textbf{Retracted} \\
\midrule
i13 & 16+ & 3 (P13, P16, excess voids) & 0 \\
i14 & 19+ & 3 (kSZ, IFMR, TDE) & 3 (time-delay) \\
i15 & 21 & 2 (P20 SNe anisotropy, P21 $D_L$ ratio) & 0 \\
i16 & 24 & 3 (P22 FRB, P23 21cm, P24 $\beta = 0$) & 0 \\
i17--i20 & 24 & 0 & 1 (GPS 59ps) \\
i21 & 26 & 2 (P25 $f \cdot \sigma_8$, P26 blue-tilted $P(k)$) & 0 \\
i22--i23 & 27 & 1 (P27 BAO amplitude) & 0 \\
i24--i25 & 29 & 2 (P28 $w_0$-$w_a$, P29 nuclear clock) & 0 \\
i26 & 31 & 2 (P30 vel.\ screening, P31 grav.\ slip) & 0 \\
i27 & 34 & 3 (P32 induced GW, P33 $A_L$, P34 phantom) & 0 \\
i28--i31 & 34 & 0 & 0 \\
\bottomrule
\end{tabular}
\end{center}

\subsection{Patent Portfolio Evolution}

\begin{center}
\begin{tabular}{lcccc}
\toprule
\textbf{Category} & \textbf{i13} & \textbf{i15} & \textbf{i19 (final)} & \textbf{Key Changes} \\
\midrule
ALIVE & 4 & 4 & 4 & P5, P7, P9, P11 \\
NICHE & 8 & 7 & 6 & P12 removed \\
THEORETICAL & 1 & 0 & 0 & P12 $\to$ DEAD \\
DEAD & 2 & 4 & 5 & +P4 (i14), +P12 (i15/i19), +P13 (i17) \\
\bottomrule
\end{tabular}
\end{center}

\subsection{Open Questions at Cycle End (i31)}

\begin{enumerate}
  \item \textbf{CMB third peak 10--61\% deficit} --- YELLOW. Only
    SymBoltz.jl can resolve.
  \item \textbf{Birefringence $\beta = 0$} --- GREEN-YELLOW. SO decisive
    2027--2028.
  \item \textbf{7 computation-only questions} --- all require SymBoltz.jl.
  \item \textbf{6 observation-only questions} --- timelines 2026--2035+.
  \item \textbf{3 computation + observation questions} --- Euclid DR1
    Oct 2026.
\end{enumerate}


\end{document}
